% -----------------------------------------------------------
\section{Natural, Nature}
% -----------------------------------------------------------
	\label{NATURAL}
	
% % ----------------------
% \subsection{See also}
% % ----------------------

% ----------------------
\subsection{1828 American Dictionary of the English Language} % *1828
% ----------------------
	\label{NATURAL-1828}

	\begin{compactitem}
		\item[1] Pertaining to nature; produced or effected by nature, or by the laws of growth, formation or motion impressed on bodies or beings by divine power. Thus we speak of the \textit{natural} growth of animals or plants; the \textit{natural} motion of a gravitating body; \textit{natural} strength or disposition; the \textit{natural} heat of the body; \textit{natural} color; \textit{natural} beauty. In this sense, \textit{natural} is opposed to artificial or acquired.
		\item[2] According to the stated course of things. Poverty and shame are the \textit{natural} consequences of certain vices.
		\item[3] Not forced; not far fetched; such as is dictated by nature. The gestures of the orator are \textit{natural}.
		\item[4] According to the life; as a \textit{natural} representation of the face. 
		\item[5] Consonant to nature. Fire and warmth go together, and so seem to carry with them as \textit{natural} an evidence as self-evident truths themselves.
		\item[6] Derived from nature, as opposed to habitual. The love of pleasure is \textit{natural} ; the love of study is usually habitual or acquired.
		\item[7] Discoverable by reason; not revealed; as \textit{natural} religion.
		\item[8] Produced or coming in the ordinary course of things, or the progress or animals and vegetables; as a \textit{natural} death; opposed to violent or premature.
		\item[9] Tender; affectionate by nature.
		\item[10] Unaffected; unassumed; according to truth and reality. What can be more \textit{natural} than the circumstances of the behavior of those women who had lost heir husbands on this fatal day?
		\item[11] Illegitimate; born out of wedlock; as a \textit{natural} son.
		\item[12] Native; vernacular; as ones \textit{natural} language.
		\item[13] Derived from the study of the works or nature; as \textit{natural} knowledge.
		\item[14] A \textit{natural} note, in music, is that which is according to the usual order of the scale; opposed to flat and sharp notes, which are called artificial.
	\end{compactitem}

% % ----------------------
% \subsection{Oxford English Dictionary} % *OED
% % ----------------------
	% \label{NATURAL-OED}

	% \begin{compactitem}
		% \item[]
	% \end{compactitem}

% ----------------------
\subsection{Scriptural Usage} % *SCRIPTURES
% ----------------------
	\label{NATURAL-SCRIPTURES}

	% % -----
	% \subsubsection{Old Testament} % *OT
	% % -----
		% \label{NATURAL-OT}
	
		% % --
		% \paragraph{KJV}
		% % --
			% \label{NATURAL-OT-KJV}
	
			% \begin{tabular}{ll}
				% Total occurrences in the OT & 
			% \end{tabular}
			
			% \begin{compactdesc}
				% \item[]
			% \end{compactdesc}
			
		% % --
		% \paragraph{Hebrew Bible}
		% % --
			% \label{NATURAL-HEBREW}
		
			% %
			% \subparagraph{\texorpdfstring{\he{sample word}}{}} % paste actual hebrew text here
			% %
				% \label{PAGA} % whatever is in step bible without periods
			
				% \begin{compactdesc}
					% \item[HALOT , PDF ] \textbf{} % Use the 1994 PDF
						% \begin{compactitem}
							% \item[1]
						% \end{compactitem}
					% \item[BDB , PDF ] \textbf{}
						% \begin{compactitem}
							% \item[1]
						% \end{compactitem}
					% \item[DCH PDF ] \textbf{}
						% \begin{compactitem}
							% \item[1]
						% \end{compactitem}
				% \end{compactdesc}
				
				% \begin{compactdesc}
					% \item[Some OT uses] \textbf{}
						% \begin{compactdesc}
							% \item[]
						% \end{compactdesc}
				% \end{compactdesc}
			
		% % --
		% \paragraph{Septuagint}
		% % --
			% \label{NATURAL-LXX}
		
			% %
			% \subparagraph{\texorpdfstring{\gr{sample word}}{}}
			% %
		
	% -----
	\subsubsection{New Testament} % *NT
	% -----
		\label{NATURAL-NT}
	
		% --
		\paragraph{KJV}
		% --
			\label{NATURAL-NT-KJV}
			
	
			\begin{tabular}{ll}
				Total occurrences in the NT & 14
			\end{tabular}
			
			\begin{compactdesc}
				\item[{\hyperref[ROMANS-1-26]{Romans 1:26}}] For this cause God gave them up unto vile affections: for even their women did change the natural use into that which is against nature:
				\item[{\hyperref[ROMANS-1-27]{Romans 1:27}}] And likewise also the men, leaving the natural use of the woman, burned in their lust one toward another; men with men working that which is unseemly, and receiving in themselves that recompence of their error which was meet.
				\item[{\hyperref[ROMANS-1-31]{Romans 1:31}}] Without understanding, covenantbreakers, without natural affection, implacable, unmerciful:
			\end{compactdesc}
			
		% --
		\paragraph{Greek New Testament} % *GREEK
		% --
			\label{NATURAL-GREEK}
		
			%
			\subparagraph{\texorpdfstring{\gr{ψυχικός}}{}}
			%
				\label{PSUCHIKOS}
			
				\begin{compactdesc}
					\item[BDAG 980b] \textbf{}
						\begin{compactitem}
							\item[1] \textbf{pertaining to the life of the natural world and whatever belongs to it, in contrast to the realm of experience whose central characteristic is \gr{πνεῦμα}, \textit{natural, unspiritual, worldly}}
							\item[1.a] adjective \gr{ψυχικὸς ἄνθρωπος} \textit{an unspiritual person}, one who merely functions bodily, without being touched by the Spirit of God (\hyperref[1CORINTHIANS-2-14]{1 Corinthians 2:14}). \gr{σῶμα ψυχικόν}, \textit{a physical body} (\hyperref[1CORINTHIANS-15-44]{1 Corinthians 15:44}). The wisdom that does not come from above is called \gr{ἐπίγειος, ψυχική} (\textit{unspiritual})\gr{, δαιμονιώδης} (\hyperref[JAMES-3-15]{James 3:15}).
							\item[1.b] substantive
							\item[1.b.$\alpha$] \gr{τὸ ψυχικόν} \textit{the physical} in contrast to \gr{τὸ πνευματικόν} (\hyperref[1CORINTHIANS-15-46]{1 Corinthians 15:46}).
							\item[1.b.$\beta$] \hyperref[JUDE-1-19]{Jude 19} calls the teachers of error \gr{ψυχικοί, πνεῦμα μὴ ἔχοντες} \textit{worldly} (literally ``psychic'') \textit{people, who do not have the Spirit}, thereby taking over the terminology of gnostic...opponents, but applying to gnostics the epithets that they used of orthodox Christians.
						\end{compactitem}
					\item[CGL 1521b, PDF 1561] \textbf{}
						\begin{compactitem}
							\item[1] relating to the soul (opposite the body); (of pleasures, achievements, or similar) \textbf{of the soul}
							\item[2] (of inclinations, boldness) \textbf{of the mind}
						\end{compactitem}
					\item[LSJ 2027b, PDF 2098] \textbf{}
						\begin{compactitem}
							\item[I1] \textit{of the soul} or \textit{life, spiritual} opposite \gr{σωματικός}
							\item[I2] \textit{of the animal life, animal} \textit{(The use in 1 Corinthians 2:14 is listed here, along with the Jude use; they gloss 1 Corinthians 2:14 as ``the \textit{natural} man.'')}
							\item[I3] \textit{brave}
							\item[II] \textit{for the soul} or \textit{spirit} of one deceased
							\item[III] \textit{cooling}
						\end{compactitem}
				\end{compactdesc}
				
				\begin{compactdesc}
					\item[Some NT uses] \textbf{}
						\begin{compactdesc}
							\item[1 Corinthians 2:14] \gr{Ψυχικὸς δὲ ἄνθρωπος οὐ δέχεται τὰ τοῦ πνεύματος τοῦ θεοῦ, μωρία γὰρ αὐτῷ ἐστίν, καὶ οὐ δύναται γνῶναι, ὅτι πνευματικῶς ἀνακρίνεται·}
							\item[1 Corinthians 15:44] \gr{σπείρεται σῶμα ψυχικόν, ἐγείρεται σῶμα πνευματικόν. Εἰ ἔστιν σῶμα ψυχικόν, ἔστιν καὶ πνευματικόν.}
							\item[1 Corinthians 15:46] \gr{ἀλλ᾽ οὐ πρῶτον τὸ πνευματικὸν ἀλλὰ τὸ ψυχικόν, ἔπειτα τὸ πνευματικόν.}
							\item[James 3:15] \gr{οὐκ ἔστιν αὕτη ἡ σοφία ἄνωθεν κατερχομένη, ἀλλὰ ἐπίγειος, ψυχική, δαιμονιώδης·}
							\item[Jude 19] \gr{οὗτοί εἰσιν οἱ ἀποδιορίζοντες, ψυχικοί, πνεῦμα μὴ ἔχοντες.}
						\end{compactdesc}
				\end{compactdesc}
		
	% % -----
	% \subsubsection{Book of Mormon} % *BOM
	% % -----
		% \label{NATURAL-BOM}
	
		% \begin{tabular}{ll}
			% Total occurrences in the BOM & 
		% \end{tabular}
		
		% \begin{compactdesc}
			% \item[]
		% \end{compactdesc}
		
	% % -----
	% \subsubsection{Doctrine and Covenants} % *DC
	% % -----
		% \label{NATURAL-DC}
	
		% \begin{tabular}{ll}
			% Total occurrences in the DC & 
		% \end{tabular}
		
		% \begin{compactdesc}
			% \item[]
		% \end{compactdesc}
		
	% % -----
	% \subsubsection{Pearl of Great Price} % *PGP
	% % -----
		% \label{NATURAL-PGP}
	
		% \begin{tabular}{ll}
			% Total occurrences in the PGP & 
		% \end{tabular}
		
		% \begin{compactdesc}
			% \item[]
		% \end{compactdesc}
		
% ----------------------
\subsection{Key Phrases} % *KEYPHRASES *PHRASES
% ----------------------
	\label{NATURAL-PHRASES}

	\begin{compactdesc}
	
		\item[natural eye] \textbf{4} | \hyperref[DC58-3]{DC 58:3}; \hyperref[MOSES-1-11]{Moses 1:11} (2); \hyperref[MOSES-6-36]{Moses 6:36}
			\begin{compactdesc}
				\item[Bible anchor verses] ---
				% \item[Commentary] 
			\end{compactdesc}
	
		\item[natural man] \textbf{6} | \hyperref[1CORINTHIANS-2-14]{1 Corinthians 2:14}; \hyperref[MOSIAH-3-19]{Mosiah 3:19} (2); \hyperref[ALMA-26-21]{Alma 26:21}; \hyperref[DC67-12]{DC 67:12}; \hyperref[MOSES-1-14]{Moses 1:14}
			\begin{compactdesc}
				\item[Bible anchor verses] \phantom{.}
					\begin{compactdesc}
						\item[1 Corinthians 2:14] \hyperref[PSUCHIKOS]{\grl{Ψυχικὸς}} \gr{δὲ ἄνθρωπος οὐ δέχεται τὰ τοῦ πνεύματος τοῦ θεοῦ, μωρία γὰρ αὐτῷ ἐστίν, καὶ οὐ δύναται γνῶναι, ὅτι πνευματικῶς ἀνακρίνεται·}
					\end{compactdesc}
				\item[Commentary] See \hyperref[NATURAL-man]{below}.
			\end{compactdesc}
			
		\item[natural mind] For this phrase, see \hyperref[MIND-PHRASES]{Mind}.
			% % \begin{compactdesc}
				% % \item[Bible anchor verses] \phantom{.}
					% % \begin{compactdesc}
						% % \item[Romans 5:5] \gr{}
					% % \end{compactdesc}
				% % \item[Commentary] ---
			% % \end{compactdesc}
			
		% \item[] \textbf{\#times} | list
			% % \begin{compactdesc}
				% % \item[Bible anchor verses] \phantom{.}
					% % \begin{compactdesc}
						% % \item[Romans 5:5] \gr{}
					% % \end{compactdesc}
				% % \item[Commentary] ---
			% % \end{compactdesc}
			
	\end{compactdesc}
	
% % ----------------------
% \subsection{Bible Dictionaries} % *DICTIONARIES
% % ----------------------
	% \label{NATURAL-BD}

% % ----------------------
% \subsection{Restoration Usage} % *RESTORATION
% % ----------------------
	% \label{NATURAL-RESTORATION}

	% % -----
	% \subsubsection{Joseph Smith} % *JS
	% % -----
		% \label{NATURAL-JS}
	
		% \begin{compactdesc}
			% \item[{\hyperref[KEY]{Date/Source}}]
		% \end{compactdesc}
		
	% % -----
	% \subsubsection{Temple Ordinances}
	% % -----
		% \label{NATURAL-TEMPLE}
	
		% \begin{compactdesc}
			% \item[{\hyperref[KEY]{Date/Source}}]
		% \end{compactdesc}
	
	% % -----
	% \subsubsection{Brigham Young} % *BY
	% % -----
		% \label{NATURAL-BY}
	
		% \begin{compactdesc}
			% \item[{\hyperref[KEY]{Date/Source}}]
		% \end{compactdesc}
	
% ----------------------
\subsection{Commentary and Studies} % *COMMENTARY *STUDIES
% ----------------------
	\label{NATURAL-COMMENTARY}

	% % -----
	% \subsubsection{Key Points}
	% % -----
		% \label{NATURAL-KEYS}
	
		% \begin{compactitem}
			% \item
		% \end{compactitem}
		
	% -----
	\subsubsection{What is the ``natural man''?}
	% -----
		\label{NATURAL-man}
		
		See my notes at \hyperref[MOSIAH-3-19]{Mosiah 3:19} and \hyperref[1CORINTHIANS-2-14]{1 Corinthians 2:14}.
		
		See BY's remarks in the sermon he gives on 19 June 1859 (a good place to find it is in the readings I prepared for the BY class in fall semester 2021; it's in ``BY Sermons 05'' on righteousness, sin, and salvation. Make sure to read LaJean's transcription, not the JD.